\documentclass[11pt]{article}
\usepackage[margin=1in]{geometry}
\usepackage{amsmath,amssymb,amsthm,mathtools,booktabs,microtype,longtable}
\usepackage[colorlinks=true,linkcolor=blue,urlcolor=blue,citecolor=blue]{hyperref}
\hypersetup{pdftitle={Sixth powers in a mean square: a guide to the 11/12 argument},
 pdfsubject={Concrete averages and the zeta zero-free-region argument}}
\allowdisplaybreaks[2]
\setlength{\emergencystretch}{2em}
\newcommand{\OO}{\mathbb Z[\omega]}
\newcommand{\N}{\mathrm N}
\newcommand{\muk}{\mu_K}
\newcommand{\ch}{\chi_{-3}}
\newcommand{\E}{\mathcal E}
\newcommand{\ind}{\mathbf 1}
\newcommand{\starred}{\sideset{}{^*}\sum}
\numberwithin{equation}{section}
\newtheorem{theorem}{Theorem}
\newtheorem{proposition}{Proposition}
\title{Sixth powers in a mean square:\\an accessible guide to the $11/12$ argument for $\zeta(s)$}
\author{Claude Opus 5.5 \quad and \quad OpenAI Astra\\[1ex]
{\normalsize at the behest of Michael O.\ Rubinstein}}
\date{}
\begin{document}
\maketitle
\begin{center}\setlength{\fboxsep}{6pt}
\fbox{\parbox{0.9\linewidth}{\small\textbf{A note from Michael Rubinstein.} These are working notes,
generated by Claude and Codex to help me read the papers on the quasi-Riemann hypothesis that
OpenAI released.}}
\end{center}
\begin{abstract}
We give an accessible account of an argument showing that $\zeta(s)$ has
no zeros in $\Re s>11/12$. The averages are written out in classical terms:
M\"obius sums over Eisenstein integers, twisted by sextic residue symbols
and averaged in mean square. Rows that are sixth powers recover the
coefficients of $1/(\zeta(s)L(s,\chi_{-3}))$, and the mean square is bounded
using Poisson summation, Kubota's cubic theta function, a quadratic large
sieve, and a descent.
\end{abstract}

\paragraph{The proof in six steps.}
\begin{enumerate}
\item To show that $\zeta(s)$ has no zeros with real part above $11/12$, it
 is enough to prove that the smoothed M\"obius sum
 $A_1(D)=\sum_n\muk(n)W(\N(n)/D)$ over Eisenstein integers, whose
 coefficients are those of $1/(\zeta(s)L(s,\chi_{-3}))$, is at most about
 $D^{11/12}$.
\item That sum is embedded in a large family of similar sums
 $A_u(D)=\sum_n\muk(n)(u/n)_6W(\N(n)/D)$, each twisted by a sextic residue
 symbol, and the roughly $H^{1/6}$ rows where $u=\pi^6$ is the sixth power
 of a prime reproduce the original sum, up to a small error.
\item The mean square $\sum_{0<\N(u)\le H}|A_u(D)|^2$ is a sum of
 nonnegative terms, so a bound of size about $DH$ for it (square-root
 cancellation in each row, on average) controls the roughly $H^{1/6}$
 copies of $|A_1(D)|^2$ it contains; this gives
 $|A_1(D)|^2\lesssim DH^{5/6}$, and taking $H$ a little larger than $D$
 gives the exponent $11/12$.
\item Poisson summation in the row variable $u$ turns this average into a
 dual family with fewer rows, and here is the key trick: the Gauss sums
 produced by Poisson summation combine with the M\"obius function to give
 \emph{cubic} Gauss sums, because the cube of a cubic Gauss sum at a prime
 carries exactly a factor $-1$, which is the M\"obius sign.
\item Cubic Gauss sums are, up to fixed residue-class factors, the Fourier
 coefficients of Kubota and Patterson's cubic theta function, whose
 symmetry under a group of matrices shortens each twisted sum and turns
 the sextic twist into an ordinary quadratic one, so that Heath-Brown's
 quadratic large sieve gives the needed bound.
\item The theta function's coefficients include extra terms indexed by
 cubes; these are removed by M\"obius inversion and by a repeated two-step
 Poisson argument that passes back through M\"obius sums and shrinks the
 number of rows each time, until the problem becomes trivial.
\end{enumerate}

\paragraph{Scope.}
This is a guide to the argument of \cite{original}, specialized to $\zeta$, not a new
independent proof of its analytic estimates. Sections~\ref{sec:coefficients}--\ref{sec:zeros}
give the deduction from one precisely stated mean-square estimate in full.
Sections~\ref{sec:poisson}--\ref{sec:descent} explain the estimates behind it,
retaining the parameters needed to see how the argument closes. The cubic
theta theorem and the quadratic large sieve are substantial inputs.
In particular, Section~\ref{sec:descent} states the uniform transfer
estimate used from \cite{original} and gives the induction that closes
once that estimate and the completed theta estimate are available.
This separates the deduction of the zero-free region from the analytic
estimates on which it rests. The finite checks in Section~\ref{sec:numerics}
test identities and parameter ranges; they do not independently verify
those uniform estimates or certify the claimed theorem.
Background on $\OO$, residue symbols, Gauss sums, Poisson summation over $\OO$
and the cubic theta function is in the companion notes \cite{QF}.

\section{The ordinary integer coefficients}\label{sec:coefficients}

Let $\mu$ be the ordinary M\"obius function and put
\begin{equation}\label{eq:chi3def}
 \ch(m)=\begin{cases}0,&3\mid m,\\1,&m\equiv1\pmod3,\\-1,&m\equiv2\pmod3.\end{cases}
\end{equation}
The coefficients of the reciprocal product are completely classical:
\begin{equation}\label{eq:cm}
 c(m)=\sum_{de=m}\mu(d)\mu(e)\ch(e),\qquad
 \frac1{\zeta(s)L(s,\ch)}=\sum_{m\ge1}\frac{c(m)}{m^s}\quad(\Re s>1).
\end{equation}
This follows by multiplying the two absolutely convergent reciprocal
Dirichlet series. For primes other than $3$, the local coefficients are
\begin{equation}\label{eq:localc}
\begin{array}{c|ccc}
 &c(p)&c(p^2)&c(p^j),\ j\ge3\\\hline
 p\equiv1\pmod3&-2&1&0\\
 p\equiv2\pmod3&0&-1&0
\end{array}
\end{equation}
and $c$ is multiplicative. Thus $c$ differs from $\mu$: for example,
$c(7)=-2$, $c(49)=1$, and $c(25)=-1$.

Choose a complex-valued smooth function
$W\in C_c^\infty((0,\infty);\mathbb C)$ supported in a fixed interval
$[a,b]\subset(0,\infty)$. Think of $W(m/D)$ as a smooth way of restricting
$m$ to size $D$. The sum to be bounded is
\begin{equation}\label{eq:A1integer}
 A_1(D)=\sum_{\substack{m\ge1\\(m,6)=1}}c(m)W(m/D).
\end{equation}
Removing the primes $2$ and $3$ gives the exact identity
\begin{equation}\label{eq:reciprocal}
 \sum_{(m,6)=1}\frac{c(m)}{m^s}
 =\frac1{(1-3^{-s})(1-4^{-s})\zeta(s)L(s,\ch)}.
\end{equation}
At $2$, the reciprocal Euler factor is $1-2^{-2s}=1-4^{-s}$;
this explains the $4$ rather than $2$ in this formula.

\section{Why the sums use Eisenstein integers}\label{sec:eisenstein}

The notes \cite[Sections~2--3]{QF} give the arithmetic of $\OO$ and its residue symbols in
full, with proofs and tables. They write $\lambda$ for $1-\omega$; their Sections~4--5
write $\sqrt{-3}$ for the element $\lambda=1+2\omega$ used here.

Put $\omega=e^{2\pi i/3}$, so $\omega^2+\omega+1=0$. An Eisenstein integer
is a point $n=a+b\omega$ with $a,b\in\mathbb Z$, and
\begin{equation}\label{eq:normdef}
 \N(n)=|n|^2=a^2-ab+b^2.
\end{equation}
These points form a triangular lattice in the plane. The ring $\OO$ has
unique factorization. Its six units are $\pm1,\pm\omega,\pm\omega^2$.
To avoid counting six associates, we choose the \emph{primary} generator:
\begin{equation}\label{eq:primarydef}
 n\equiv1\pmod3,\quad\text{equivalently }a\equiv1\pmod3,\ b\equiv0\pmod3.
\end{equation}
Every element coprime to $3$ has exactly one primary associate. In sums
below, $n,f,b_1$ and divisors of primary elements are primary and coprime
to $6$, unless stated otherwise. The row variables $u,k,h$ run over all
of $\OO$. A star means that the element is squarefree.

Define $\muk(n)=0$ if a prime square divides $n$, and
$\muk(n)=(-1)^r$ if $n$ is a product of $r$ distinct primary primes;
$\muk(1)=1$. This is the same factorization rule as for the ordinary
M\"obius function, applied in $\OO$.

A rational prime $p\equiv1\pmod3$ splits into two conjugate prime
factors of norm $p$. A prime $p\equiv2\pmod3$ remains prime, of norm
$p^2$. Consequently
\begin{equation}\label{eq:normgroup}
 \sum_{\substack{n\ \mathrm{primary}\\\N(n)=m}}\muk(n)=c(m)
 \qquad((m,6)=1).
\end{equation}
For a split prime the local factor is $(1-p^{-s})^2$; for an inert
prime it is $1-p^{-2s}$. These are exactly the factors in
\eqref{eq:cm}. Thus \eqref{eq:A1integer} can also be written
\begin{equation}\label{eq:A1eis}
 A_1(D)=\sum_n\muk(n)W(\N(n)/D).
\end{equation}
The two-dimensional notation keeps track of the two primes above a split
rational prime separately. The norm alone loses this information.

\section{The average, written out}\label{sec:average}

For a prime $\pi$ coprime to $6$, the number $\N(\pi)-1$ is divisible
by $6$. Define the sextic residue symbol by
\begin{equation}\label{eq:sexticdef}
 \left(\frac{u}{\pi}\right)_6
 =\text{the sixth root of unity congruent to }
 u^{(\N(\pi)-1)/6}\pmod\pi
\end{equation}
when $\pi\nmid u$, and define it to be $0$ when $\pi\mid u$.
Extend it multiplicatively in the denominator and write
$\chi_n(u)=(u/n)_6$. It is multiplicative in the numerator as well.
For $n=1$ its value is $1$, including at $u=0$.
Whenever we take a power of a character, its zero values remain zero,
also for negative or zero exponents.

The individual sums, and their mean square, are
\begin{align}
 A_u(D)&=\sum_n\muk(n)\chi_n(u)W(\N(n)/D),\label{eq:Au}\\
 M(D,H)&=\sum_{0<\N(u)\le H}|A_u(D)|^2.\label{eq:M}
\end{align}
There are about $H$ rows $u$, and each row contains about $D$ possible
columns $n$. More explicitly, with $Q(a,b)=a^2-ab+b^2$,
\begin{equation}\label{eq:coordinates}
 M(D,H)=\sum_{\substack{x,y\in\mathbb Z\\0<Q(x,y)\le H}}
 \left|\sum_{\substack{a\equiv1\ (3),\ b\equiv0\ (3)\\(a+b\omega,2)=1}}
 \muk(a+b\omega)
 \left(\frac{x+y\omega}{a+b\omega}\right)_6
 W\left(\frac{Q(a,b)}D\right)\right|^2.
\end{equation}
This is an ordinary sum over integer coordinates. Coprimality to $2$
means that $a,b$ are not both even.

Expanding the square gives the requested double sum over the columns:
\begin{equation}\label{eq:double}
 \boxed{\ M(D,H)=\sum_{n_1,n_2}\muk(n_1)\muk(n_2)
 \overline{W(\N(n_1)/D)}W(\N(n_2)/D)S_H(n_1,n_2),\ }
\end{equation}
where the finite character correlation is
\begin{equation}\label{eq:S}
 S_H(n_1,n_2)=\sum_{0<\N(u)\le H}
 \overline{\left(\frac{u}{n_1}\right)_6}
 \left(\frac{u}{n_2}\right)_6.
\end{equation}
Every coefficient in this average is now specified by factorization,
finite-field exponentiation, and a smooth weight.

For completeness, one can group every row by ordinary integer norms:
\begin{equation}\label{eq:cu}
 c_u(m)=\sum_{\substack{n\ \mathrm{primary}\\\N(n)=m}}
 \muk(n)\chi_n(u),\qquad
 A_u(D)=\sum_{(m,6)=1}c_u(m)W(m/D).
\end{equation}
Here $c_1(m)=c(m)$ from \eqref{eq:cm}. In general $c_u(m)$ requires the
residue symbols: $\mu$ and $\ch$ alone do not specify the separate twists
at conjugate primes.

\paragraph{The analytic estimate to prove.}
For every fixed $0<\vartheta\le1/10$ and every $\varepsilon>0$, the
required estimate is
\begin{equation}\label{eq:MS}
 \boxed{\quad M(D,D^{1+\vartheta})
 \ll_{W,\vartheta,\varepsilon}D^{2+\vartheta+\varepsilon}.
 \quad}
\end{equation}
The notation $\ll$ means an inequality up to a constant independent of
$D$. The factor $D^\varepsilon$ allows an arbitrarily small power loss.
The full argument controls the constant using finitely many derivatives
of $W$. The estimate says that the \emph{average} squared row sum is
at most $D^{1+\varepsilon}$; it does not assert such a bound for every row.

\section{The diagonal and the special rows}\label{sec:diagonal}

\subsection{Why \texorpdfstring{$HD$}{HD} is the natural scale}
For $n_1=n_2=n$, the correlation is simply
\begin{equation}\label{eq:diagcount}
 S_H(n,n)=\#\{u\in\OO:0<\N(u)\le H,\ (u,n)=1\}.
\end{equation}
For fixed $n$, elementary lattice counting and inclusion--exclusion give
\begin{equation}\label{eq:diagasym}
 S_H(n,n)=\frac{2\pi}{\sqrt3}H
 \prod_{\pi\mid n}\left(1-\frac1{\N(\pi)}\right)
 +O\bigl(\tau(n)(\sqrt H+1)\bigr),
\end{equation}
where $\tau(n)$ is the number of primary divisors of $n$. The area of the
disc is $\pi H$ and the lattice has covolume $\sqrt3/2$.
For the proof, the simpler uniform upper bound $S_H(n,n)\ll H$ suffices.
The diagonal contribution to \eqref{eq:double} is therefore $O_W(HD)$.
Estimate \eqref{eq:MS} bounds the whole mean square on this scale, up to
$D^\varepsilon$. It controls the aggregate off-diagonal contribution;
it does not bound the sum of the absolute values of individual
off-diagonal terms.

\subsection{Sixth powers expose the untwisted sum}
If $u=v^6$, then
\begin{equation}\label{eq:sixthpower}
 \chi_n(v^6)=\ind_{(v,n)=1},\qquad
 A_{v^6}(D)=\sum_{(n,v)=1}\muk(n)W(\N(n)/D).
\end{equation}
We use prime $v=\pi$ of comparable size so that the omitted terms are few.
Put $Y=H^{1/6}$ and
\begin{equation}\label{eq:PY}
 \mathcal P_Y=\{\pi\text{ primary prime},\ \pi\nmid6:
                         Y/2<\N(\pi)\le Y\},\qquad J=|\mathcal P_Y|.
\end{equation}
Then $\N(\pi^6)\le H$, and
\begin{equation}\label{eq:primeerror}
 |A_{\pi^6}(D)-A_1(D)|\ll_W D/Y.
\end{equation}
Indeed, a contributing omitted $n$ is $\pi m$ with
$\N(m)\le bD/\N(\pi)$; a disc of squared radius $T$ contains $O(1+T)$
lattice points. In our range $D/Y\ge1$.
The prime number theorem in the progression $1\pmod3$ gives
$J\asymp Y/\log Y$. These rows are distinct: equal sixth powers have
associated bases, and primary associates are unique.

The elementary inequality $|z|^2\le2|w|^2+2|z-w|^2$, followed by
positivity of \eqref{eq:M}, gives the rigorous extraction
\begin{equation}\label{eq:extract}
 J|A_1(D)|^2
 \le2\sum_{\pi\in\mathcal P_Y}|A_{\pi^6}(D)|^2
       +O_W(JD^2/Y^2)
 \le2M(D,H)+O_W(JD^2/Y^2).
\end{equation}
The dyadic interval $Y/2<\N(\pi)\le Y$ is useful: it makes the error
$D/Y$ uniform for every selected row.

\subsection{Cubes, squares, and rational rows}
There are two other useful exact identities:
\begin{equation}\label{eq:cubesquare}
 \chi_n(v^3)=\left(\frac vn\right)_2,\qquad
 \chi_n(v^2)=\left(\frac vn\right)_3.
\end{equation}
Thus cube rows are quadratic twists and square rows are cubic twists.
These classes overlap, as sixth powers illustrate.
For a nonzero \emph{rational integer} $v$ and an odd denominator $n$,
\begin{equation}\label{eq:rationalquad}
 \left(\frac vn\right)_2=\left(\frac v{\N(n)}\right),
\end{equation}
with the ordinary Jacobi symbol on the right. For split primes this is
Euler's criterion over $\mathbb F_p$; for an inert prime of norm $p^2$,
every nonzero element of $\mathbb F_p$ is a square in $\mathbb F_{p^2}$.
Both sides vanish if $p\mid v$. Consequently
\begin{equation}\label{eq:cuberational}
 A_{v^3}(D)=\sum_{(m,6)=1}c(m)\left(\frac vm\right)W(m/D).
\end{equation}
If $\chi_v(m)=(v/m)$ on odd $m$ denotes the corresponding quadratic
Dirichlet character there, the Dirichlet series of these coefficients is
\begin{equation}\label{eq:cubeeuler}
 \prod_{p\nmid6}(1-\chi_v(p)p^{-s})
                    (1-\chi_v(p)\ch(p)p^{-s}).
\end{equation}
This states the Euler factors explicitly without conventions at $2$.

A rational numerator does not generally eliminate the need for
Eisenstein residue symbols in a sextic or cubic twist. Also, the rational
slice $u\in\mathbb Z$, $0<\N(u)\le H$, has $\asymp H^{1/2}$ rows, since
$\N(u)=u^2$. The whole lattice has $\asymp H$ rows. Sixth-power lattice
rows number $\asymp H^{1/6}$; sixth powers with rational bases number
$\asymp H^{1/12}$. These are different counts.

\section{From the mean square to the zero-free region}\label{sec:zeros}

\begin{theorem}
If \eqref{eq:MS} holds for every smooth compactly supported $W$ and every
fixed $0<\vartheta\le1/10$, then $\zeta(s)$ has no zeros in $\Re s>11/12$.
\end{theorem}
\begin{proof}
Take $H=D^{1+\vartheta}$ in \eqref{eq:extract}. Absorbing a logarithm into
$D^\varepsilon$ gives
\begin{align}
 |A_1(D)|^2
 &\ll_{W,\vartheta,\varepsilon}DH^{5/6}D^\varepsilon+D^2H^{-1/3}\label{eq:A1square}\\
 &=D^{11/6+5\vartheta/6+\varepsilon}
   +D^{5/3-\vartheta/3}.\label{eq:A1exponents}
\end{align}
Taking square roots, and then choosing $\vartheta$ and $\varepsilon$
sufficiently small, yields, for every $\eta>0$,
\begin{equation}\label{eq:mobiusbound}
 A_1(D)\ll_{W,\eta}D^{11/12+\eta}.
\end{equation}
The exponent $11/12$ is half of $2-1/6$: the mean square has size roughly
$D^2$, and the useful rows number roughly $D^{1/6}$. Quadratic or cubic twists would have more
useful rows and a better exponent; Section~\ref{sec:whysextic} explains why the argument needs
sextic ones.

Here is the precise link with zeros. Write
\begin{equation}\label{eq:Fdef}
 F(s)=(1-3^{-s})(1-4^{-s})\zeta(s)L(s,\ch),\qquad
 \widehat W(s)=\int_0^\infty W(y)y^s\,\frac{dy}{y}.
\end{equation}
Suppose $\zeta(\rho)=0$ and $\Re\rho>11/12$. Choose a nonzero
nonnegative smooth function $\phi$ supported in $(1,2)$ and set
$W(y)=y^{-\rho}\phi(y)$. Then $\widehat W(\rho)>0$.
By \eqref{eq:mobiusbound},
\begin{equation}\label{eq:BW}
 B_W(s)=\int_0^\infty A_1(D)D^{-s}\,\frac{dD}{D}
\end{equation}
is holomorphic in $\Re s>11/12$. There is no difficulty at $D=0$:
the sum vanishes for $D<1/2$. For $\Re s>1$, absolute convergence and
\eqref{eq:reciprocal} give $F(s)B_W(s)=\widehat W(s)$.
Meromorphic continuation extends this identity to $\Re s>11/12$.
Since $L(s,\ch)$ is entire and $\rho\ne1$, $F(\rho)=0$.
Evaluation at $\rho$ now gives $0=\widehat W(\rho)>0$, a contradiction.
\end{proof}

The remaining question is how to prove \eqref{eq:MS}. The next three
sections explain the argument of \cite{original}, specialized to $\zeta$, and identify its hard inputs.

\section{Poisson summation changes the coefficients}\label{sec:poisson}

Additive characters, Gauss sums and Poisson summation over $\OO$, including the identities
used in this section, are explained in \cite[Section~4]{QF}.

\subsection{A smooth row cutoff}
Replace the disc cutoff by a fixed nonnegative smooth radial weight
$\Phi(\N(u)/H)$ with $\Phi(t)\ge1$ for $0\le t\le1$. One can choose
$\Phi$ rapidly decreasing and its two-dimensional Fourier transform
$\widehat\Phi$ compactly supported. This gives an upper bound for $M$.
It is this smoothed average to which Poisson summation is applied.

Here is an exact version of the transformation. Put
\begin{equation}\label{eq:poissonnotation}
 \lambda=1+2\omega=i\sqrt3,\qquad
 e(z)=\exp(4\pi i\,\operatorname{Im}z/\sqrt3),\qquad
 \widehat\Phi(\N(w))=\frac2{\sqrt3}
     \int_{\mathbb C}\Phi(|z|^2)e(-zw)\,dx\,dy.
\end{equation}
Only squarefree $n_1,n_2$ occur because of the M\"obius factors. Write
$n_i=gz_i$ with $g=(n_1,n_2)$. Then $z_1,z_2$ are coprime and
\begin{equation}\label{eq:pairedchar}
 \overline{\chi_{n_1}(u)}\chi_{n_2}(u)
 =\ind_{(u,g)=1}\psi(u),\qquad
 \psi=\overline{\chi_{z_1}}\chi_{z_2}.
\end{equation}
The character $\psi$ is primitive modulo $z_1z_2$. For a primitive
character modulo $q$, set
\begin{equation}\label{eq:gammapsi}
 \gamma(\psi)=\N(q)^{-1/2}\sum_{x\bmod q}\psi(x)e(x/q).
\end{equation}
Inclusion--exclusion for $(u,g)=1$ and lattice Poisson summation give
\begin{equation}\label{eq:poissonexact}
 \begin{split}
 &\sum_{u\in\OO}\Phi(\N(u)/H)\ind_{(u,g)=1}\psi(u)\\
 &\quad=\frac{H\gamma(\psi)}{\sqrt{\N(z_1z_2)}}
 \sum_{d\mid g}\frac{\muk(d)\psi(d)}{\N(d)}
 \sum_{h\in\OO}\overline{\psi(h)}
 \widehat\Phi\left(\frac{H\N(h)}{\N(d)\N(z_1z_2)}\right).
 \end{split}
\end{equation}
The zero frequency $h=0$ occurs only when $z_1=z_2=1$, exactly the
diagonal. All other frequencies satisfy
\begin{equation}\label{eq:freqrange}
 \N(h)\ll\frac{\N(d)\N(z_1z_2)}H
       \ll\frac{D^2}{H\N(g)},
\end{equation}
because $d\mid g$ and $\N(z_i)\asymp D/\N(g)$. Thus a row range of size
$H>D$ becomes a frequency range of size at most a constant times
$D^2/H<D$. The divisor $d$ in the exact formula must be retained.

\subsection{The M\"obius--Gauss identity}
For squarefree $n$ put
\begin{equation}\label{eq:gammaj}
 \gamma_j(n)=\N(n)^{-1/2}\sum_{x\bmod n}\chi_n(x)^j e(x/n),\qquad
 \alpha(n)=\frac n{|n|},\qquad G(n)=\overline{\chi_n(4)}\gamma_3(n).
\end{equation}
The arithmetic identity that makes the method work is
\begin{equation}\label{eq:gauss}
 \muk(n)\gamma_{-1}(n)
 =\chi_n(-1)G(n)^{-1}\overline{\alpha(n)}\gamma_2(n).
\end{equation}
The right side contains a \emph{cubic} Gauss sum, since
$\chi_n^2=(\cdot/n)_3$. The unit-modulus factor $G$ depends only on
the residue class of $n$ modulo $72$ (in fact modulo $4$ here).

To see its effect on \eqref{eq:double}, let $\mathfrak G$ be the finite
group of invertible classes modulo $72$ that are $1$ modulo $3$.
Expand the function $t\mapsto G(t^{-1})$ in its finite group characters:
$G(t^{-1})=\sum_\xi c_\xi\xi(t)$. With
$a_\xi(n)=\overline{\alpha(n)}\gamma_2(n)\xi(n)$, reciprocity and
\eqref{eq:gauss} give, for coprime squarefree $z_1,z_2$,
\begin{equation}\label{eq:paired}
 \muk(z_1)\muk(z_2)\gamma(\overline{\chi_{z_1}}\chi_{z_2})
 =\sum_\xi c_\xi a_\xi(z_1)\overline{a_\xi(z_2)}.
\end{equation}
There are only finitely many $\xi$, independently of $D$. Formula
\eqref{eq:paired} preserves the form of a mean square while replacing
its M\"obius coefficients by cubic Gauss coefficients.

After resolving the coprimality conditions and dividing the variables
into ranges $F\le\N(f)<2F$, the averages to be bounded are
\begin{equation}\label{eq:energy}
 \E(\mathcal H,X,F)=\frac1{XF}
 \starred_{F\le\N(f)<2F}\ \sum_{0<\N(k)\le\mathcal H}
 \left|\starred_n a_\xi(n)\chi_n(k)\chi_n(f)^4W(\N(n)/X)\right|^2.
\end{equation}
Here $\xi,W$ are suppressed from the notation. Both $f,n$ are squarefree.
The fourth power in $\chi_n(f)^4$ has two sources. Writing $e$ for the
exclusion divisor in Poisson summation gives a twist $\chi_n(he^5)$.
Removing the coprimality condition between the two columns introduces
a common divisor $v$ and the identity
$a_\xi(vm)=a_\xi(v)a_\xi(m)\chi_m(v)^4$. With $f=ev$ and $k=eh$,
\begin{equation}\label{eq:fourthpower}
 he^5v^4=k(ev)^4=kf^4.
\end{equation}
Thus the $v$-part comes from twisted multiplicativity, and the $e$-part
from the Poisson exclusion divisor. If $B\ge1$ is a further divisor-size
parameter, the initial ranges are
\begin{equation}\label{eq:initial}
 X=\frac D{BF},\qquad \Sigma:=XF=\frac DB,\qquad
 \mathcal H\ll\frac{D^2}{HB^2}.
\end{equation}
The needed bound is $\E\ll D^\varepsilon\Sigma$. The normalizations and
divisor sums then recover \eqref{eq:MS}; this is the first Poisson
reduction.

\section{Why cubic theta and a quadratic large sieve appear}\label{sec:theta}

\subsection{Adding the cube indices}
The cubic Gauss coefficients in \eqref{eq:energy} occur, after a fixed
normalization, in the Fourier expansion of Kubota's cubic theta function.
Its indices include cube factors, powers of $\lambda$, and units.
The explicit coefficient formulas are due to Patterson and are recorded
in \cite[Section 5.1 and Appendix]{DR}. The notes \cite[Section~5]{QF} explain the theta
function and state the properties used here.

Here is the normalization needed for our sum. Restrict a Fourier index
$\ell$ to those for which $\lambda^3\ell$ is primary and coprime to $6$.
Among the nonzero coefficients, this selects exactly
$\ell=\lambda^{-3}nb_1^3$, with $n$ squarefree and $b_1$ arbitrary,
both primary and coprime to $6$. The coefficient of the conjugated theta
function at such an index is
\begin{equation}\label{eq:thetacoef}
 d_0(\lambda^{-3}nb_1^3)
 =3^{5/2}|b_1|\,\widetilde g(n)
 =3^{5/2}|b_1|\,\overline{\left(\frac\lambda n\right)_3}\gamma_2(n).
\end{equation}
Eisenstein's supplementary law says that $(\lambda/n)_3$ depends only
on $n\bmod9$. Multiplying the twist at $x=\lambda^3\ell=nb_1^3$ by
$(\lambda/x)_3=(\lambda/n)_3$ cancels this extra factor. It is absorbed
into the same finite congruence-class bookkeeping used above.
This selection applies to the original series; unit factors and varying
powers of $\lambda$ reappear in the reflected cusp expansions and are
included in their estimates.

For fixed $k,f$, put
$\Psi(n)=\xi(n)\chi_n(k)\chi_n(f)^4$. The convenient completed sum is
\begin{equation}\label{eq:T}
 T(X;k,f)=\frac1{\sqrt X}\starred_n\sum_{b_1}
 \overline{\alpha(n)}\gamma_2(n)\Psi(n)
 \overline{\alpha(b_1)}^3\Psi(b_1)^3\N(b_1)^{1/2}
 W\left(\frac{\N(n)\N(b_1)^3}X\right).
\end{equation}
The $b_1=1$ part is precisely $X^{-1/2}$ times the column sum in
\eqref{eq:energy}. All sums here are finite because $W$ has compact support.

\subsection{The transformation and its purpose}
The theta function obeys a transformation law relating its values at
rational translates to its values near one of three cusps. Expanding
$\Psi$ in finite additive characters makes that law applicable to
\eqref{eq:T}. Differentiating in the horizontal variable removes the
constant Fourier terms; Mellin transformation then gives a reflected sum.
The constant terms would otherwise produce residues, so their removal
is an essential step.

The important local calculation is especially simple to state.
Let $\pi\nmid6$ be a prime at which the original twist is $\chi_\pi^j$, and expand
$\chi_\pi^j$ in additive characters modulo $\pi$. Characters are extended by
zero at multiples of their prime, also when an exponent is $0$.

If $1\le j\le5$, the zero frequency has coefficient $0$, so every frequency
with a nonzero coefficient is nonzero modulo $\pi$, and $\pi$ divides the
denominator of every translate. The reflected coefficient
then contains, up to a factor of absolute value $1$,
\begin{equation}\label{eq:Blocal}
 B_{\pi,j}(x)=
 \begin{cases}
  \chi_\pi(x)^{-j-2},&j\notin\{0,4\},\\
  \N(\pi)^{-1/2}(-1+\N(\pi)\ind_{\pi\mid x}),&j=4.
 \end{cases}
\end{equation}

If $j=0$, the twist is the principal character $\chi_\pi^0(x)=\ind_{\pi\nmid x}$, and its
expansion has two parts:
\begin{equation}\label{eq:jzero}
 \ind_{\pi\nmid x}=\Bigl(1-\frac1{\N(\pi)}\Bigr)
 -\frac1{\N(\pi)}\sum_{h\not\equiv0\ (\pi)}e\Bigl(\frac{hx}\pi\Bigr).
\end{equation}
The constant part (the \emph{inactive} branch) gives translates whose denominators
are not divisible by $\pi$. It contributes only the factor $1-\N(\pi)^{-1}$, with no
reflected local factor at $\pi$. The terms with $h\not\equiv0$ (the \emph{active}
branch) give, up to a factor of absolute value $1$,
\begin{equation}\label{eq:jzeroactive}
 B_{\pi,0}(x)=\N(\pi)^{-1/2}\chi_\pi(x)^{-2}.
\end{equation}
Both branches occur. In particular,
\begin{equation}\label{eq:Bone}
 B_{\pi,1}(x)=\chi_\pi(x)^3=\left(\frac x\pi\right)_2.
\end{equation}
Thus the squarefree part of the row twist becomes \emph{quadratic}.
If $c$ is the denominator of the rational translate in lowest terms, put
$Q=\N(c)$. The reflected weight has argument $\N(\ell)X/Q^2$, so the
effective dual norm range is $\N(\ell)\lesssim Q^2/X$. It is shorter
than the original range $X$ when $Q<X$. Rapid decay becomes useful when
$X/Q^2$ is large: the nonzero frequencies $\ell\in\lambda^{-4}\OO$ have $\N(\ell)\ge1/81$,
so their weight arguments then grow uniformly. The inequality $Q^2<X$
alone is not a quantitative negligibility estimate.

The quadratic large sieve says, in the form needed here,
\begin{equation}\label{eq:qls}
 \starred_{\substack{\N(k)\le K\\(k,6)=1}}
 \left|\starred_{\N(n)\le U}\beta(n)\left(\frac nk\right)_2\right|^2
 \ll_\varepsilon(KU)^\varepsilon(K+U)\starred_n|\beta(n)|^2.
\end{equation}
Here the generators are primary, the column $n$ may be even, and the
coefficients $\beta(n)$ are arbitrary. Both the row and column squarefree
restrictions are essential for this bound. In particular, if square
columns $n=b_1^2$ were allowed, then $(b_1^2/k)_2=1$ whenever
$(b_1,k)=1$, so coefficients supported on squares would reinforce each
other instead of cancelling. This is the substantial large-sieve input; it is a specialization of the
number-field quadratic large sieve of \cite{GL}.

Applying \eqref{eq:qls} to the reflected sum, accounting for the local
factors at primes dividing $f$, including both active and inactive branches,
and then treating the square factors of
$k$, gives the completed estimate
\begin{equation}\label{eq:completed}
 \boxed{\quad\sum_{0<\N(k)\le\mathcal H}|T(X;k,f)|^2
 \ll D^\varepsilon\left(\mathcal H+
                   \frac{\mathcal H^2\N(f)}X\right).\quad}
\end{equation}
Parameters are bounded by fixed powers of $D$, and a fixed finite
smoothness norm of $W$ is understood. The estimate is derived from the theta transformation and the large sieve,
as in \cite[Proposition~5.2 and Section~6]{original}.

For clarity, the passage to all rows uses $k=usv^2$, with $u$ a unit and
$s$ squarefree. The identity
$\chi_n(v^2)\chi_n(f)^4=\chi_n(fv^2)^4$ absorbs the square into $f$.
The corresponding squarefree-row estimate allows this more general
second argument. Summing over $v$ costs $\sum_v\N(v)^{-2}<\infty$.

\subsection{Why the twists are sextic}\label{sec:whysextic}
Twists of any order $k$ could be used to amplify $A_1(D)$. For $k=2,3,6$, replace
$\chi_n(u)$ in \eqref{eq:Au} by $\chi_n(u)^{6/k}$, which is $(u/n)_2$, $(u/n)_3$ or $(u/n)_6$ by
\eqref{eq:cubesquare}. The rows $u=v^k$ then reproduce the untwisted sum, as in \eqref{eq:sixthpower},
and there are about $H^{1/k}$ of them. A mean-square bound of the diagonal size $HD$ would give,
exactly as in Section~\ref{sec:zeros} with $H\approx D$,
\begin{equation}\label{eq:orderk}
 |A_1(D)|^2\ll\frac{HD}{H^{1/k}}\approx D^{2-1/k},\qquad\text{so}\qquad A_1(D)\ll D^{1-1/(2k)+\varepsilon}.
\end{equation}
Quadratic twists would give $3/4$, cubic twists $5/6$, and sextic twists only $11/12$. So the order $6$
is chosen for the available proof of the mean-square bound. The difference
lies in what Poisson summation in $u$ does to the M\"obius coefficients. As in \eqref{eq:paired}, the
coefficient of a column $z$ in the dual sum is $\muk(z)$ times the Gauss sum of the order-$k$ character
$\chi_z^{\pm6/k}$.
\begin{itemize}
\item \emph{Quadratic twists.} Poisson summation produces the quadratic Gauss sums $\gamma_3(n)$. These
are explicit: by \eqref{eq:gammaj}, $\gamma_3(n)=\chi_n(4)G(n)$. Here $G$ depends only on $n$ modulo
$72$, and $\chi_n(4)=(2/n)_3$ depends only on $n$ modulo $2$ by cubic reciprocity, so $\gamma_3(n)$ depends
only on $n$ modulo $72$. So the dual coefficients are still $\muk(n)$, up to a periodic phase, and the dual sum is again a
mean square of M\"obius sums with quadratic twists. Its square frequencies $k=v^2$ carry untwisted
M\"obius sums once more, so this transformation by itself supplies no new estimate. The quadratic large sieve
\eqref{eq:qls} cannot be applied to the original family either: it requires squarefree rows, and the
rows $u=v^2$ that carry $A_1(D)$ are exactly the ones it must exclude.
\item \emph{Cubic twists.} Poisson summation produces the cubic Gauss sums $\gamma_{\pm2}(n)$. The
identities behind \eqref{eq:gauss} give
\begin{equation}\label{eq:cubicdual}
 \muk(n)\gamma_2(n)=\chi_n(-1)G(n)\,\alpha(n)\gamma_{-1}(n)
\end{equation}
for squarefree $n$. The M\"obius function is absorbed, but into a \emph{sextic} Gauss sum.
This points toward theta functions on a sixfold metaplectic cover.
Such theta functions do have automorphy, and Kubota's Dirichlet series
built from higher-order Gauss sums have functional equations \cite{BB}.
For the sextic theta function, Br\"oker and Hoffstein \cite{BH} prove a formula for the
coefficients at $\pi^4$ for certain primes $\pi$, and conjecture the coefficients at $\pi^2$
and the squares of those at $\pi$ from numerical evidence; the coefficients at squarefree
indices are not known.
What is missing from the present argument is a coefficient description
and a uniform estimate for the transformed sums that play the roles
of \eqref{eq:thetacoef} and \eqref{eq:completed}.
The failure of Whittaker uniqueness in this higher-cover setting helps
explain the difficulty \cite{KP,BH}; it does not imply that transformation
formulas do not exist.
\item \emph{Sextic twists.} Poisson summation produces the sextic Gauss sums $\gamma_{-1}(n)$, and
\eqref{eq:gauss} turns $\muk(n)\gamma_{-1}(n)$ into $\overline{\alpha(n)}\gamma_2(n)$, up to a periodic
phase. The reason is the factorization, from the Hasse--Davenport product formula,
\begin{equation}\label{eq:hdfactor}
 \gamma_1(n)=\overline{\chi_n(4)}\,\gamma_2(n)^2\gamma_3(n),\qquad \gamma_2(n)^3=\muk(n)\alpha(n).
\end{equation}
A sextic Gauss sum is the square of a cubic one times a quadratic one. The quadratic factor is
explicit, and the cube of the cubic factor supplies exactly the M\"obius sign. The resulting coefficient
is a Fourier coefficient of the cubic theta function, which Patterson determined. Its automorphy then
turns the remaining sextic twist into a quadratic one, \eqref{eq:Bone}. Only at that stage, when the
coefficients are Gauss sums rather than values of the M\"obius function, is the quadratic large sieve applied.
\end{itemize}
Of these three routes, the sextic one supplies the coefficient conversion
and transformed estimates used here. The price is the count $H^{1/6}$ of
amplifying rows, which gives $11/12$ rather than $3/4$ or $5/6$.
Notice that the order of the twists and the degree of the auxiliary
metaplectic cover are different: sextic twists lead to the \emph{cubic}
theta function, whereas the cubic-twist calculation above produces
\emph{sextic} Gauss sums.

Better control of higher metaplectic theta coefficients could therefore
help improve this method, but it is one possible route, not a necessary
gateway to a larger zero-free region. Different averages or amplification
could also help. Formula~\eqref{eq:orderk} shows that simply increasing
the order of the twists makes this particular amplification weaker:
perfect-power rows become rarer.

\section{Removing the added cubes and closing the argument}\label{sec:descent}

The purpose of descent is to bound the original average
$\E(\mathcal H,X,F)$ in \eqref{eq:energy}. In plain terms, the completed
theta estimate handles small cube indices, while the remaining large
cube indices lead to averages of exactly the same kind with fewer rows.
The point requiring care is that all the other conditions, including
the smooth weights, must survive this reduction.

One cannot obtain a bound for the $b_1=1$ part of \eqref{eq:T} simply
by discarding the other terms: cancellation may occur inside $T$.
Instead use the exact M\"obius inversion
\begin{equation}\label{eq:inversion}
 \frac1{\sqrt X}\starred_n a_\xi(n)\chi_n(k)\chi_n(f)^4W(\N(n)/X)
 =\sum_h\frac{\muk(h)\overline{\alpha(h)}^3\Psi(h)^3}{\N(h)}
           T(X/\N(h)^3;k,f).
\end{equation}
In this formula $h$ is primary and coprime to $6$. Substituting
\eqref{eq:T} and collecting the total cube index uses
$\sum_{h\mid b_1}\muk(h)=\ind_{b_1=1}$, proving the identity.

\subsection{Small cubes are controlled by the completed estimate}
Choose
\begin{equation}\label{eq:Hc}
 H_c^3=\min\left(X,\frac{X^2}{\mathcal H^2}\right),\qquad
 L_{b_1}=X/\N(b_1)^3.
\end{equation}
For $\N(h)\le H_c$, estimate \eqref{eq:completed} at scale
$X/\N(h)^3$ gives terms bounded by
\begin{equation}\label{eq:smallcubes}
 \mathcal H+\frac{\mathcal H^2F\N(h)^3}X\ll\Sigma,
 \qquad\Sigma=XF,
\end{equation}
provided $\mathcal H\le\Sigma$. Cauchy--Schwarz costs only divisor sums
and logarithms, which fit into $D^\varepsilon$.
Expanding the remaining completed sums in \eqref{eq:inversion} leaves
the same original average at smaller column scales. The resulting
cube-removal estimate \cite[Lemma~5.3]{original} is
\begin{equation}\label{eq:cuberemoval}
 \E(\mathcal H,X,F)\ll D^\varepsilon
 \left(\Sigma+
   \sup_{\substack{\N(b_1)>H_c\\L_{b_1}>1}}
                  \E(\mathcal H,L_{b_1},F)\right).
\end{equation}
Scales $L_{b_1}\le1$ are bounded by counting. If $\mathcal H^2\le X$,
the supremum is empty and the desired estimate already follows.
Here and below, the suppressed dependence on $W$ is through a fixed
finite smoothness norm. More precisely, the $\Sigma$ term in
\eqref{eq:cuberemoval} is $\Sigma\|W\|_{C^{J_c}(I)}^2$ for a finite
integer $J_c$, with the same $W$ in each average in the supremum.

\subsection{Large cubes force a smaller row range}
Keep $\Sigma=XF$ from the original average fixed while replacing the
column length $X$ by $L=L_{b_1}$. In particular, do not reset $\Sigma$
to $LF$. The extra room $LF\le\Sigma$ is what permits the enlargement
below.
For the remaining terms, two more Poisson summations exchange cubic
Gauss coefficients for M\"obius coefficients and then back again.
The paired identity \eqref{eq:paired} is used in both directions.
The first summation gives frequencies of norm about $L^2/\mathcal H$,
with additional divisor factors. Combining the frequency with the
auxiliary variable as $y=hf^2$ gives a natural M\"obius-row range of size
$L^2F^2/\mathcal H$. A crucial step is to enlarge this range by positivity
before the second Poisson summation. Suppressing the common divisor
factor $\N(d)/\N(C)^2$ and fixed support constants, the two scales are
\begin{equation}\label{eq:Yscales}
 Y_{\rm natural}=\frac{L^2F^2}{\mathcal H},\qquad
 Y_{\rm enlarged}=\frac{\Sigma LF}{\mathcal H},\qquad
 \frac{Y_{\rm enlarged}}{Y_{\rm natural}}=\frac{\Sigma}{LF}\ge1.
\end{equation}
After separating the smooth weights, the row sums are nonnegative mean
squares, so the larger smooth majorant gives a valid upper bound.
The enlarged range makes the second Poisson frequency range smaller.
For a remaining divisor $r$ in a range $R\le\N(r)<2R$, it gives
\begin{equation}\label{eq:Hprime}
 \mathcal H'=\frac{\mathcal HL}{\Sigma FR^2}.
\end{equation}
Without enlargement the corresponding scale would be
$\mathcal H/(F^2R^2)$, which need not decrease when $F=R=1$.

After the second summation, regrouping the divisor terms \emph{before}
taking absolute values forces a divisibility condition $f'\mid k'$.
At each prime the cancellation is
\begin{equation}\label{eq:localcancel}
 (1+\ind_{\pi\mid k'})-\ind_{\pi\mid k'}
                         -\ind_{\pi\nmid k'}=\ind_{\pi\mid k'}.
\end{equation}
This divisibility keeps the next average in the ranges needed for
induction. The kernel after the second summation vanishes unless
$0<\N(k')\le\mathcal HL/(\Sigma F\N(r)^2)$. Retain only those $r$ for
which a nonzero term occurs. If $F'\le\N(f')<2F'$, then $f'\mid k'$ gives
\begin{equation}\label{eq:retained}
 F'\le\N(f')\le\N(k')\le\frac{\mathcal HL}{\Sigma F\N(r)^2}\le\mathcal H',
\end{equation}
so $\mathcal H'\ge F'\ge1$. It also keeps the column scales at least $1$.
Put $X'=L/(RF')$. Removing an exclusion divisor $j\mid r$ produces the
replacement
\begin{equation}\label{eq:exclusionswap}
 (X',F')\longmapsto\bigl(X'/\N(j),\ F'\N(j)\bigr).
\end{equation}
It preserves both $\mathcal H'$ and the product $\Sigma'=X'F'=L/R$,
and hence the ratio $\mathcal H'/\Sigma'$. Since $\N(j)\le\N(r)$,
$R\le\N(r)$, $F\ge1$ and
$\mathcal H\le\Sigma$, \eqref{eq:retained} gives
\begin{equation}\label{eq:childcolumn}
 \frac{X'}{\N(j)}\ge\frac{\Sigma F\N(r)^2}{\mathcal HR\N(j)}
 \ge\frac{\Sigma F}{\mathcal H}\ge1.
\end{equation}

\subsection{The precise transfer estimate}
We now state the analytic input needed from \cite{original}, in the
form sufficient for this guide. Write $\E(\mathcal H,X,F;\xi,W)$ when
the twist and weight in \eqref{eq:energy} need to be shown. For an interval
$I$, put $\|W\|_{C^m(I)}=\max_{0\le r\le m}\sup_{x\in I}|W^{(r)}(x)|$.

\begin{proposition}[Transfer estimate {\cite[Proposition~5.4]{original}}]\label{prop:guide-transfer}
Fix $C_0\ge1$, $D\ge2$, a compact interval $I\subset(0,\infty)$,
and $W\in C_c^\infty(I)$. Suppose
$1\le\mathcal H,L,F,\Sigma\le D^{C_0}$, with
$\max\{\mathcal H,LF\}\le\Sigma$. Here $\Sigma$ is an independent
target scale. For every integer $m\ge0$ and $\delta>0$,
\begin{equation}\label{eq:transferprop}
 \E(\mathcal H,L,F;\xi,W)
 \ll D^\delta\Sigma\|W\|_{C^{4m+12}(I)}^2
 \left(1+\sup\frac{\E(\mathcal H',X',F';\xi',U)}{\Sigma'}\right).
\end{equation}
The supremum runs over averages of the same form, with
$\mathcal H',X',F'\ge1$, $\Sigma'=X'F'$, and
\begin{equation}\label{eq:transfer}
 \mathcal H'\le\frac{\mathcal HL}{\Sigma F},\qquad
 \frac{\mathcal H'}{\Sigma'}\le\frac{\mathcal H}{\Sigma},\qquad
 \Sigma'\le L.
\end{equation}
For $I=[u,v]\subset(0,\infty)$, the new weights satisfy
$U\in C_c^\infty([u/16,4v])$ and $\|U\|_{C^m}\le1$.
The characters $\xi'$ belong to the same fixed finite family of
residue-class twists. An empty supremum is zero. The implied constant
may depend on $m,C_0,\delta,I$ and the fixed arithmetic data, but not
on $D$ or the moving scales.
\end{proposition}

The source proves the stronger version with a smooth majorant of the
row sum on the left. Its proof, in \cite[Section~7]{original}, is the two-Poisson argument just
outlined, including the bounds for the separated smooth kernels and
zero frequencies. The full uniform estimate is a substantive analytic
input: the three scale inequalities \eqref{eq:transfer}, by themselves,
do not establish it.

\subsection{Why the induction closes}
This subsection follows the proof of \cite[Proposition~5.1]{original}, given in
\cite[Section~5.5]{original}.
Fix $\kappa>0$, $C_0\ge1$, and $\varepsilon>0$, and assume
\begin{equation}\label{eq:standing}
 1\le\mathcal H,X,F,\qquad \Sigma=XF\le D^{C_0},\qquad
 \mathcal H\le\Sigma D^{-\kappa}.
\end{equation}
The inductive assertion, for each integer $j\ge0$, is
\begin{equation}\label{eq:inductive}
 \E(\mathcal H,X,F;\xi,W)
 \ll D^\varepsilon\Sigma\|W\|_{C^{J_j}(I)}^2
 \qquad\text{whenever }\mathcal H\le D^{j\kappa}.
\end{equation}
This assertion is made for every $\varepsilon>0$ and every compact
interval $I\subset(0,\infty)$. The finite derivative order $J_j$ may
depend on $j,\kappa,C_0,\varepsilon$, but not on $D$ or $I$; the implied
constant may also depend on $I$. This formulation allows the support
interval to enlarge at a transfer step.

\paragraph{Base case.}
If $j=0$, then $\mathcal H\le1$, so there are at most six nonzero
rows. There are $O(F)$ choices of $f$ and $O_I(X)$ column terms, each
of absolute value at most $\|W\|_\infty$. The normalization in
\eqref{eq:energy} therefore gives
\begin{equation}\label{eq:basecase}
 \E(\mathcal H,X,F;\xi,W)
 \ll_I\frac{F}{XF}(X\|W\|_\infty)^2
 =X\|W\|_\infty^2\le\Sigma\|W\|_\infty^2.
\end{equation}

\paragraph{Induction step.}
Assume the assertion at level $j$, and let
$\mathcal H\le D^{(j+1)\kappa}$. Use that assertion with loss
$\varepsilon/3$, and let $m$ be its required derivative order.
Apply cube removal \eqref{eq:cuberemoval} and
Proposition~\ref{prop:guide-transfer}, each with loss $\varepsilon/3$.
If the cube-removal supremum is empty, its $\Sigma$ term already
suffices. Otherwise $\N(b_1)^3>X^2/\mathcal H^2$.
For each remaining $L=L_{b_1}$ we have $LF\le\Sigma$, so the transfer
estimate applies with the unchanged target $\Sigma=XF$. Every child
average satisfies
\begin{equation}\label{eq:contraction}
 \mathcal H'\le\frac{\mathcal H}{\N(b_1)^3F^2}
 <\mathcal H\left(\frac{\mathcal H}{XF}\right)^2
 \le\mathcal H D^{-2\kappa}.
\end{equation}
In particular, $\mathcal H'\le D^{j\kappa}$. The other hypotheses survive
as well:
\begin{equation}\label{eq:survive}
 \frac{\mathcal H'}{\Sigma'}\le\frac{\mathcal H}{\Sigma}
 \le D^{-\kappa},\qquad
 \Sigma'\le L\le\Sigma\le D^{C_0},\qquad
 \mathcal H',X',F'\ge1.
\end{equation}
The induction hypothesis applies to the child weight $U$ on
$[u/16,4v]$. Since $\|U\|_{C^m}\le1$, it gives the uniform bound
\begin{equation}\label{eq:childbound}
 \sup\frac{\E(\mathcal H',X',F';\xi',U)}{\Sigma'}
 \ll D^{\varepsilon/3}.
\end{equation}
The transfer contributes a second factor $D^{\varepsilon/3}$ and
cube removal a third. Taking $J_{j+1}\ge\max\{4m+12,J_c\}$ proves the
assertion at level $j+1$. Thus both the derivative requirement and
the total $D^\varepsilon$ loss are controlled.

Finally take $j=\lceil C_0/\kappa\rceil$; this covers every allowed
$\mathcal H$. The number of induction levels depends on $C_0,\kappa$,
not on $D$. Suppressing the resulting fixed smoothness norm again,
we obtain
\begin{equation}\label{eq:finalenergy}
 \E(\mathcal H,X,F)\ll D^\varepsilon\Sigma
 \qquad(\mathcal H\le\Sigma D^{-\kappa}).
\end{equation}
No estimate for zeros of $\zeta$ is assumed in this induction: its
inputs are cube removal and the transfer estimate, and its starting
case is counting. Establishing those analytic inputs is the work that
must precede this formal closure.

Finally, the initial scales \eqref{eq:initial}, with $H=D^{1+\vartheta}$,
satisfy
\begin{equation}\label{eq:initialratio}
 \frac{\mathcal H}{\Sigma}\ll\frac{D^{-\vartheta}}B.
\end{equation}
For sufficiently large $D$ this is at most $D^{-\vartheta/2}$.
Thus \eqref{eq:finalenergy} applies with $\kappa=\vartheta/2$.
Empty ranges and bounded column scales are handled by counting.
The first Poisson reduction gives \eqref{eq:MS}, and
Section~\ref{sec:zeros} then gives the zero-free half-plane.

\paragraph{The roles of the ingredients.}
The ordinary convolution $c=\mu*(\mu\ch)$ identifies the reciprocal
zeta product. Sixth-power rows recover that convolution inside a larger
nonnegative average. Poisson summation shortens the row range and exposes
cubic Gauss sums. The theta transformation converts their twists into
quadratic ones, to which a large sieve applies. Cube inversion and the
contraction \eqref{eq:contraction} transfer the resulting completed bound
back to the original sums.

\section{Numerical checks of Sections \ref{sec:coefficients}--\ref{sec:descent}}\label{sec:numerics}

The identities and parameter inequalities listed below, from Sections~\ref{sec:coefficients}--\ref{sec:descent},
numbered or not, were tested on finite samples. The checks use Python with exact arithmetic in $\OO$, written as pairs
$a+b\omega$, and the conventions of these notes: primary means $\equiv1\pmod3$,
$\lambda=1+2\omega$, $\chi_n(x)=(x/n)_6$, and $e(z)=\exp(4\pi i\operatorname{Im}z/\sqrt3)$. Residue
symbols are computed by Euler's criterion modulo each prime, and Gauss sums by brute force over
complete residue systems. In Sections~\ref{sec:coefficients} and~\ref{sec:eisenstein}, $c(m)$ is computed
from its definition~\eqref{eq:cm} with the ordinary M\"obius function, and $\muk(n)$ from the
factorization of $\N(n)$ in $\mathbb Z$, so that the two sides of~\eqref{eq:normgroup} are computed
independently. Sums of complex numbers are in double precision, and the Mellin and Bessel
integrals use $25$--$30$-digit arithmetic. ``Exact'' below means that the two sides agreed
identically, or to rounding error, in every case tried. The scripts, with a description of each check,
are at \url{https://github.com/mrubinst/L-function_zero_free_strips/tree/main/zeta_11_12_tests}.

\begin{center}
\small
\renewcommand{\arraystretch}{1.15}
\begin{longtable}{@{}p{2.3cm}p{8.6cm}p{3.6cm}@{}}
\toprule
Display & What was checked & Result\\\midrule
\endhead
\multicolumn{3}{@{}l}{\emph{Section~\ref{sec:coefficients}}}\\
\eqref{eq:chi3def} & $\ch(m)$ equals the Kronecker symbol $(-3/m)$ for $m\le3000$ and is completely
 multiplicative ($2000$ random pairs) & exact\\
\eqref{eq:cm} & coefficientwise, $c*(1*\ch)=\delta$ for all $m\le10^5$, i.e.\ the series of $c$ times
 $\zeta(s)L(s,\ch)$ is $1$; the truncated series $\sum_{m\le10^5}c(m)m^{-s}$ against
 $1/(\zeta(s)L(s,\ch))$ at $s=3$, $3+4i$, $4$ & exact; $10^{-13}$\\
\eqref{eq:localc} & the table for all $9690$ prime powers $p^j\le10^5$ with $p\ne3$, and $c(3)=-1$,
 $c(3^j)=0$ for $j\ge2$; $c(m)=\prod_{p^e\Vert m}c(p^e)$ for all $m\le10^5$; the values $c(7)$,
 $c(49)$, $c(25)$ & exact\\
\eqref{eq:reciprocal} & coefficientwise for all $m\le10^5$; the local factors at $2$ and $3$ are
 $1-4^{-s}$ and $1-3^{-s}$; numerically at $s=3+2i$ & exact; $7\cdot10^{-14}$\\
\multicolumn{3}{@{}l}{\emph{Section~\ref{sec:eisenstein}}}\\
\eqref{eq:normdef} & $\omega^2+\omega+1=0$; $|a+b\omega|^2=a^2-ab+b^2$ for $|a|,|b|\le40$; there are
 exactly six units, $\pm1,\pm\omega,\pm\omega^2$ & exact\\
\eqref{eq:primarydef} & $n\equiv1\pmod3$ in $\OO$ if and only if $a\equiv1$, $b\equiv0\pmod3$; each
 $n$ coprime to $3$ has exactly one primary associate, and a multiple of $\lambda$ has none
 ($|a|,|b|\le60$) & exact\\
splitting & for every prime $p<3000$: if $p\equiv1\pmod3$, there are exactly two primary elements of
 norm $p$, they are conjugate, and their product is $p$; if $p\equiv2\pmod3$, no element has norm $p$
 and $-p$ is primary; $3=-\lambda^2$ & exact\\
\eqref{eq:normgroup} & the number of primary $n$ with $\N(n)=m$ equals $\sum_{d\mid m}\ch(d)$ for all
 $m\le10^5$ with $3\nmid m$ (unique factorization); $\sum_{\N(n)=m}\muk(n)=c(m)$ for all $33\,333$
 values $m\le10^5$ coprime to $6$; the local factors $(1-p^{-s})^2$ and $1-p^{-2s}$ & exact\\
\eqref{eq:A1integer}, \eqref{eq:A1eis} & $A_1(D)$ computed both ways, with a complex smooth $W$,
 for $D=300$, $3000$, $40\,000$ & $10^{-14}$\\
\multicolumn{3}{@{}l}{\emph{Section~\ref{sec:average}}}\\
\eqref{eq:sexticdef} & $6\mid\N(\pi)-1$ for the first $3000$ primary primes $\pi\nmid6$; $\chi_n$ is
 multiplicative in the numerator and in the denominator (random samples) & exact\\
\eqref{eq:coordinates} & the coordinate conditions $a\equiv1$, $b\equiv0\pmod3$, $(a,b)$ not both even
 are equivalent to ``primary and coprime to $6$'' ($|a|,|b|\le40$); the coordinate sum equals
 $M(D,H)$ for $D=60$, $H=90$ and a complex weight & exact\\
\eqref{eq:double}, \eqref{eq:S} & the double sum over columns, with $S_H$ computed directly, equals $M(D,H)$
 & $5\cdot10^{-16}$\\
\eqref{eq:cu} & $A_u(D)=\sum_{(m,6)=1}c_u(m)W(m/D)$ for $40$ rows; $c_1(m)=c(m)$ for all $m\le2000$; an
 example showing that $c_u(p)$ depends on the residue symbols & $10^{-15}$; exact\\
\eqref{eq:MS} & sanity check of the size: $M(D,D^{1.1})/(D^{2.1})=0.35,\ 0.32,\ 0.31$ for
 $D=100,200,400$ & consistent\\
\multicolumn{3}{@{}l}{\emph{Section~\ref{sec:diagonal}}}\\
\eqref{eq:diagcount}, \eqref{eq:diagasym} & lattice counts for four $n$ and $H\le10^4$:
 $|\text{count}-\text{main term}|\le0.23\,\tau(n)(\sqrt H+1)$ & consistent\\
\eqref{eq:sixthpower} & $\chi_n(v^6)=\ind_{(v,n)=1}$ (random $v,n$); $A_{v^6}(D)$ for prime and composite $v$
 & exact; $10^{-14}$\\
\eqref{eq:primeerror} & $A_{\pi^6}-A_1=-\sum_{\pi\mid n}\muk(n)W$ for eight $\pi$, $D=1500$; the ratio to
 $D/\N(\pi)$ stays bounded & $10^{-16}$\\
\eqref{eq:PY} & $J\log Y/Y=0.50,\ 0.52,\ 0.51$ for $Y=10^3,10^4,10^5$; the rows $\pi^6$ are distinct
 & consistent\\
\eqref{eq:extract} & $J|A_1|^2\le2\sum|A_{\pi^6}|^2+2\sum|A_{\pi^6}-A_1|^2$ on actual values ($D=1500$,
 $Y=200$, $J=21$); the error term is $0.036\,JD^2/Y^2$ & holds\\
\eqref{eq:cubesquare}, \eqref{eq:Bone} & $\chi_n(v^3)$ and $\chi_n(v^2)$ against the quadratic and cubic
 symbols computed independently by Euler's criterion & exact\\
\eqref{eq:rationalquad} & $(v/n)_2=(v/\N(n))$ for $|v|\le30$ and $150$ values of $n$ & exact\\
\eqref{eq:cuberational} & $c_{v^3}(m)=c(m)(v/m)$ for all $m\le2000$ and $v\in\{2,-1,3,5,-3,7,-6,10\}$
 & exact\\
\eqref{eq:cubeeuler} & the Euler product has coefficients $c(m)(v/m)$, $m\le2000$ & exact\\
\multicolumn{3}{@{}l}{\emph{Section~\ref{sec:zeros}}}\\
\eqref{eq:A1square}--\eqref{eq:mobiusbound} & the exponent algebra & symbolic\\
\eqref{eq:reciprocal}, \eqref{eq:Fdef} & $\sum_{(m,6)=1}c(m)m^{-s}=1/F(s)$ at $s=3$ and $s=2.5+4i$
 & within truncation error\\
\eqref{eq:BW} & the substitution $\int W(m/D)D^{-s}\,dD/D=m^{-s}\widehat W(s)$ & $10^{-35}$\\
\multicolumn{3}{@{}l}{\emph{Section~\ref{sec:poisson}}}\\
\eqref{eq:poissonnotation} & $\widehat\Phi$ for $\Phi(t)=e^{-\pi t}$ by quadrature, against
 $(2/\sqrt3)e^{-4\pi t/3}$ & $10^{-21}$\\
\eqref{eq:pairedchar} & the factorization of $\overline{\chi_{n_1}}\chi_{n_2}$; $\psi$ is primitive
 ($128$ local components) & exact\\
\eqref{eq:poissonexact} & the exact Poisson formula with $\Phi(t)=e^{-\pi t}$, five cases including
 $g\ne1$ and $z_1=1$ & $5\cdot10^{-15}$\\
\eqref{eq:double}+\eqref{eq:poissonexact} & end to end: the whole smoothed mean square, computed directly,
 equals its Poisson-dual form ($D=40$, $H=60$) & $2\cdot10^{-15}$\\
\eqref{eq:freqrange}, \eqref{eq:initial}, \eqref{eq:initialratio} & the scale algebra & symbolic\\
\eqref{eq:gammaj}, \eqref{eq:gauss} & the M\"obius--Gauss identity for $70$ squarefree $n$, including
 composite $n$ & $5\cdot10^{-15}$\\
\eqref{eq:gammaj} & $G(n)$ depends only on $n\bmod4$ & exact\\
\eqref{eq:paired} & $22$ coprime pairs; a control with the conjugation on the wrong factor fails
 (error $1.3$) & $4\cdot10^{-15}$\\
\eqref{eq:fourthpower} & $a_\xi(vm)=a_\xi(v)a_\xi(m)\chi_m(v)^4$ ($22$ pairs) and the character identity
 behind $he^5v^4=kf^4$ & $4\cdot10^{-15}$\\
\multicolumn{3}{@{}l}{\emph{Section~\ref{sec:theta}}}\\
\eqref{eq:thetacoef} & against the coefficient table of the theta function as transcribed from~\cite{DR},
 $104$ indices; $(\lambda/n)_3$ depends only on $n\bmod 9$; the selection of indices
 $\lambda^{-3}nb_1^3$ ($908$ indices) & $4\cdot10^{-15}$; exact\\
theta input & automorphy under $z\mapsto z+1$ and the inversion & $2\cdot10^{-16}$\\
Section~\ref{sec:theta} & the transformation with one $\bar z$-derivative, end to end at a single cusp:
 the twisted coefficient sum against the reflected sum with the kernel
 $\prod_\pm\Gamma(1+t\pm\frac16)/\Gamma(1-t\pm\frac16)$, for $39$ cusp configurations (all cusp
 types, all three values of the multiplier). Controls with the sign flipped, with the kernel
 $\Gamma(\frac43\pm t)\Gamma(\frac53\pm t)$, or with the wrong power of $\overline{\alpha(c)}$ fail at
 size $1$ & $5.5\cdot10^{-9}$\\
\eqref{eq:Blocal}--\eqref{eq:jzeroactive} & the additive expansions of $\chi_\pi^j$ ($1\le j\le5$) and
 of $\ind_{\pi\nmid x}$, and the reflected Gauss sums
 $\sum_{h\ne0}\chi_\pi(h)^{-j-2}e(\varepsilon h^{-1}x/\pi)$, for eight primes; the active
 $j=0$ branch, summed directly for several $\varepsilon$, equals
 $-\gamma_2(\pi)\chi_\pi(\varepsilon)^{-2}\N(\pi)^{-1/2}\chi_\pi(x)^{-2}$, so it has modulus
 $\N(\pi)^{-1/2}$ when $\pi\nmid x$ and vanishes when $\pi\mid x$ & $3\cdot10^{-14}$\\
text after \eqref{eq:Bone} & $\lambda^4=9$, so $\N(\ell)=\N(m)/81$ for $\ell=\lambda^{-4}m$; over
 $|a|,|b|\le30$ the least nonzero value is $1/81$, attained exactly at the six units divided by $9$
 & exact\\
\eqref{eq:qls} & sanity check with $K=U=400$ ($109$ squarefree rows, $138$ squarefree columns,
 even columns included): the ratio of the left side to $(K+U)\sum|\beta|^2$ is at most $0.254$ for six
 coefficient choices & consistent\\
\eqref{eq:completed} & sanity check: $\sum_k|T(X;k,f)|^2\approx(0.2\text{--}0.35)\mathcal H$ for $X=1500$,
 $\mathcal H\le1000$, $f\in\{1,7+9\omega\}$ & consistent\\
text before \eqref{eq:cubicdual} & $\gamma_3(n)$ depends only on $n\bmod4$ ($90$ squarefree $n$, $12$ classes),
 and $\chi_n(4)$ only on $n\bmod2$ & exact\\
\eqref{eq:cubicdual}, \eqref{eq:hdfactor} & for $70$ squarefree $n$, including composite $n$ & $6\cdot10^{-15}$\\
\multicolumn{3}{@{}l}{\emph{Section~\ref{sec:descent}}}\\
\eqref{eq:inversion} & the cube inversion for three pairs $(k,f)$ at $X=4000$, with six nontrivial
 cubes & $2\cdot10^{-15}$\\
\eqref{eq:Hc}--\eqref{eq:retained}, \eqref{eq:childcolumn}, \eqref{eq:transfer}, \eqref{eq:contraction}
 & the inequalities between scales, on $200\,000$ random admissible
 parameter sets, and \eqref{eq:retained}, \eqref{eq:childcolumn} on a further $200\,000$ sets with a
 retained $r$. A control that keeps only $F'\le\mathcal H'$ instead of \eqref{eq:retained} violates
 $X'/\N(j)\ge\Sigma F/\mathcal H$ in about $9\,000$ sets. \eqref{eq:localcancel} exactly & no violation\\
\eqref{eq:exclusionswap} & the inclusion--exclusion identity behind the replacement, for the column sums
 (\cite[Lemma~4.4]{original}): $8$ cases at $X=3000$, with $r$ having one to three prime factors and two
 characters $\xi$. Also the resulting inequality
 $\E_r(\mathcal H,X,F)\le\tau(r)\sum_{d\mid r}\E(\mathcal H,X/\N(d),F\N(d))$, both sides computed from
 \eqref{eq:energy}, for three parameter sets & $3\cdot10^{-15}$; holds\\
\eqref{eq:exclusionswap}, \eqref{eq:childcolumn} & the replacement preserves $\mathcal H'$, $\Sigma'=L/R$
 and $\mathcal H'/\Sigma'$, and the new scales are at least $1$ ($200\,000$ random sets) & no violation\\
\eqref{eq:basecase} & on computed values of \eqref{eq:energy} with $\mathcal H\le1$, four pairs $(X,F)$ & holds\\
\eqref{eq:transfer}, \eqref{eq:standing}, \eqref{eq:contraction}, \eqref{eq:survive} & $100\,000$ random
 chains of induction steps (cube removal, transfer with a retained $r$, exclusion removal). At every step the
 hypotheses \eqref{eq:standing} and $\max\{\mathcal H,LF\}\le\Sigma$ hold with $\Sigma=XF$ unchanged, and
 $\mathcal H'\le D^{(j-1)\kappa}$; every chain ends within $\lceil C_0/\kappa\rceil$ steps & no violation\\
\eqref{eq:transferprop}, \eqref{eq:finalenergy} & sanity check: $\E/(\Sigma\|W\|_\infty^2)\le0.16$ for
 $X=300$--$1200$, $F\in\{1,7\}$, $\mathcal H\le160$, falling like $1/X$ at fixed $\mathcal H$. With
 $a(n)=1$ instead, the row $k=1$ has no cancellation and alone gives about $0.012\,\Sigma\|W\|_\infty^2$
 & consistent\\
\bottomrule
\end{longtable}
\end{center}

These checks confirm the identities and scale algebra on the samples tried.
They do not prove the uniform analytic estimates \eqref{eq:MS},
\eqref{eq:qls}, \eqref{eq:completed}, \eqref{eq:cuberemoval}, or
Proposition~\ref{prop:guide-transfer}. Finite sanity checks cannot
establish such bounds. Section~\ref{sec:descent} proves the implication
from its stated analytic inputs to \eqref{eq:finalenergy}; the
numerical checks do not replace the proofs of those inputs.

\begin{thebibliography}{9}
\bibitem{original} OpenAI, \emph{The Quasi-Riemann Hypothesis},
manuscript dated October 5, 2026.
\href{https://github.com/openai/math/tree/main/preprints/The-Quasi-Riemann-Hypothesis-October-5-2026}{Manuscript and source on GitHub}.
\bibitem{DR} A. Dunn and M. Radziwi\l\l,
\emph{Bias in cubic Gauss sums: Patterson's conjecture},
Annals of Mathematics (2) \textbf{200} (2024), 967--1057.
\url{https://arxiv.org/abs/2109.07463}.
\bibitem{QF} \emph{Quadratic number fields, with $\mathbb Q(\sqrt{-3})$ worked out in detail},
companion notes, 2026.
\url{https://www.math.uwaterloo.ca/~mrubinst/quasi_rh_notes/quadratic_fields_notes.pdf};
\href{https://github.com/mrubinst/L-function_zero_free_strips/tree/main/quadratic_fields_notes}{source on GitHub}.
\bibitem{KP} D. A. Kazhdan and S. J. Patterson, \emph{Metaplectic forms},
Publications Math\'ematiques de l'IH\'ES \textbf{59} (1984), 35--142.
\bibitem{GL} L. Goldmakher and B. Louvel,
\emph{A quadratic large sieve inequality over number fields},
Mathematical Proceedings of the Cambridge Philosophical Society
\textbf{154} (2013), 193--212.
\url{https://arxiv.org/abs/1112.1642}.
\bibitem{BB} B. Brubaker and D. Bump,
\emph{On Kubota's Dirichlet series},
Journal f\"ur die reine und angewandte Mathematik \textbf{598} (2006), 159--184.
\href{https://doi.org/10.1515/CRELLE.2006.073}{doi:10.1515/CRELLE.2006.073}.
\bibitem{BH} R. Br\"oker and J. Hoffstein,
\emph{Fourier coefficients of sextic theta series}, preprint, 2013.
\url{https://arxiv.org/abs/1312.0568}.
\end{thebibliography}
\end{document}
